\documentclass[11pt]{article}
\usepackage[margin=2.1cm]{geometry}
\usepackage{booktabs}
\usepackage{amsmath}
\usepackage{enumitem}
\usepackage{xcolor}
\usepackage[T1]{fontenc}
\usepackage{lmodern}

\setlength{\parindent}{0pt}
\setlength{\parskip}{3.6pt}
\newcommand{\code}[1]{\texttt{#1}}
\newcommand{\hr}{\vspace{2pt}\hrule\vspace{6pt}}

\begin{document}

\begin{center}
{\Large\bfseries EE\,362 --- Computer Organization Laboratory}\\[2pt]
{\large\bfseries Lab 10 --- Instruction Decoder I}\\[4pt]
Fall 2026 \quad$\cdot$\quad Lab TA: Gandhi Torres \quad$\cdot$\quad Sections: Mon / Tue, 3:45--5:45~PM, ECIII 237
\end{center}
\hr

\section*{1. Read this first: the scope is small}

This CPU implements \textbf{eight instructions} --- the ones its program uses.
Not forty. You do four this week and four next week.

\textbf{The instruction set manual is not your specification. The control table
is}, and it is given to you below, filled in. You implement it; you do not have
to derive it. Deriving it is the lecture's job, not fifty-five minutes at a
bench.

\section*{2. Objective}

\begin{quote}
\emph{A control unit is one truth table with sixteen output columns.}
\end{quote}

In Lab 1 you took a four-row truth table and wrote one Boolean equation from it.
Today you take an eight-row table and write sixteen.

\textbf{That is the entire difference.} Nothing new has been introduced since
week one. The control unit of a processor \emph{is} a truth table --- that is
not an analogy, it is literally what you are about to write.

\section*{3. Background}

\subsection*{Instruction fields}

\begin{center}
\begin{tabular}{llllll}
\toprule
R-type & \code{opcode} 6 & \code{rs} 5 & \code{rt} 5 & \code{rd} 5 \quad \code{sa} 5 & \code{func} 6 \\
I-type & \code{opcode} 6 & \code{rs} 5 & \code{rt} 5 & \multicolumn{2}{l}{\code{immediate} 16} \\
J-type & \code{opcode} 6 & \multicolumn{4}{l}{\code{index} 26} \\
\bottomrule
\end{tabular}
\end{center}

The opcode is \emph{always} at bits 31--26. It has to be: until you have decoded
it you do not know which format you are looking at, so the field you need first
must sit at a fixed position. Likewise \code{rs} and \code{rt} are in the same
place across R and I types, which lets the register file start reading before
decoding finishes. \textbf{Field placement in an ISA is a hardware decision.}

\subsection*{The control bundle}

Your decoder produces \textbf{sixteen} control signals. Sixteen separate ports
through every entity would be unreadable and would guarantee a wiring error, so
they are bundled into one VHDL \code{record}, \code{CTRL\_REC}, provided
complete in \code{CPU\_Pkg.vhd}. In hardware nothing changes --- a record is
sixteen wires either way --- but at this size it is the difference between a
design you can debug and one you cannot.

\subsection*{The default, and why it is a safety property}

\code{CTRL\_DEFAULT} sets the safe values: no register write, no memory access,
no bus driving, PC advances sequentially. \textbf{Assign it first, then
override only what each instruction changes.}

Skip it and the synthesiser must keep unassigned fields unchanged --- which
means \textbf{sixteen latches}, and control signals that depend on which
instruction ran \emph{before}. Example 3 demonstrates this: the same \code{LUI}
decodes differently depending on whether an \code{ADDI} preceded it.

In Lab 12 the symptom is the worst kind: every instruction works when tested
alone, the program still fails, and it fails \emph{differently} depending on
instruction order.

You have met this exact bug twice already --- Lab 1's missing \code{else} (where
you wanted the latch) and Lab 3's missing \code{when others} (where you did
not). Same sentence every time: \emph{no assignment means hold the previous
value.}

\section*{4. This week's four rows}

Only the \emph{differences} from the default are shown. ``--'' means leave the
default.

\begin{center}
\small
\begin{tabular}{llllllllll}
\toprule
instr & opcode/func & ALU\_func & ALU\_Asel & ALU\_Ssel & IR\_Xsel & IR\_Xlsft & RF\_Dsel & RF\_strb & RF\_Wsel \\
\midrule
\code{ADD}  & \code{000000}/\code{100000} & ADD & AOUT & -- & -- & -- & ALU & STROBE & RD \\
\code{SLL}  & \code{000000}/\code{000000} & SLL & AOUT & SA & -- & -- & ALU & STROBE & RD \\
\code{ADDI} & \code{001000} & ADD & IMMED & -- & SIGN & NONE & ALU & STROBE & RT \\
\code{LUI}  & \code{001111} & ADD & IMMED & -- & ZERO & NONE & ALU & STROBE & RT \\
\bottomrule
\end{tabular}
\end{center}

\subsection*{Three things worth noticing in that table}

\begin{itemize}[leftmargin=1.6em,itemsep=2pt]
\item \code{ADD} and \code{SLL} write \code{rd}; \code{ADDI} and \code{LUI}
write \code{rt}. \textbf{R-type and I-type put the destination in different
fields.} That is what \code{RF\_Wsel} is for, and getting it backwards writes
the wrong register \emph{silently}.

\item \code{ADDI} sign-extends and \code{LUI} zero-extends. Same field, same
width, different meaning --- decided by the instruction. This is the Lab 4
question arriving for real: get it wrong and \code{ADDI r1,r1,-1} adds 65535
instead of subtracting one, and the Lab 12 loop counts the wrong way forever.

\item \code{ADD} and \code{SLL} never read the extender output, so
\code{IR\_Xsel} and \code{IR\_Xlsft} are \textbf{don't cares} for them. Leaving
them at the default is correct --- and next week you will see why that matters.
\end{itemize}

\section*{5. Guided examples}

\begin{description}[leftmargin=2.4em,style=nextline,itemsep=2pt]
\item[Example 1 --- \code{fields}] All three formats extracted at once, on real
instructions.
\item[Example 2 --- \code{decode\_add}] One instruction decoded end to end. Five
signals set, eleven defaults kept. This is the template.
\item[Example 3 --- \code{latch\_storm}] Two decoders differing in one line. One
assert is \emph{expected} to fail.
\end{description}

\section*{6. Assignment}

Complete \code{Decoder.vhd}: the four rows above. Structure is a \code{case} on
the opcode with a nested test on \code{func} for the \code{SPECIAL} family ---
\textbf{\code{SPECIAL} is not one instruction, it is a family}, so both fields
must be checked.

\subsection*{Testbench}

\textbf{At least one vector per instruction}, plus all of: an unimplemented
opcode (inert); \code{SPECIAL} with an unimplemented \code{func} (also inert);
and \textbf{an order-dependence test} --- feed \code{ADDI} then \code{LUI} and
check \code{LUI} still zero-extends, then the reverse.

That last one is the check a single-instruction-at-a-time testbench can never
make, and it is the one that catches the worst bug in this lab.

\section*{7. Evidence to collect (E1, E2, E4, E7)}

\begin{center}
\begin{tabular}{lll}
\toprule
& Source & What it must show \\
\midrule
E1 & Flow Summary + Messages & 0 errors and \textbf{no inferred-latch warning} \\
E2 & RTL Viewer & combinational logic, \emph{no storage} \\
E4 & Resource Utilisation & \textbf{zero registers} \\
E7 & Transcript & every control signal asserted, 0 failures \\
\bottomrule
\end{tabular}
\end{center}

\textbf{Any register or latch reported means you have the Example 3 bug.}

\section*{8. Report --- \code{lab10.tex}}

\begin{itemize}[leftmargin=1.6em,itemsep=1pt]
\item \textbf{Specification:} your four rows of the control table.
\item \textbf{Design:} why the default must be assigned first; why
\code{SPECIAL} needs the \code{func} field; and one sentence on any signal whose
value surprised you.
\item \textbf{Verification:} the order-dependence test, and why it is necessary.
\item \textbf{Discussion:} required.
\end{itemize}

\section*{9. Grading (100 points)}

\begin{center}
\begin{tabular}{lr}
\toprule
Pre-lab, submitted at the start of the session & 5 \\
Design compiles cleanly, warnings explained & 25 \\
Testbench and test vectors, coverage justified & 20 \\
Simulation results match the specification & 20 \\
Report chapter quality & 15 \\
In-lab demonstration and explanation & 15 \\
\midrule
Total & 100 \\
\bottomrule
\end{tabular}
\end{center}

Late work is accepted for one week at $-20$ points.

\section*{10. Things that will go wrong}

\begin{itemize}[leftmargin=1.6em,itemsep=1pt]
\item \textbf{No default assignment.} Sixteen latches. This will happen to
several people; the Messages panel will say so, and the register count will not
be zero.
\item \textbf{\code{RF\_Wsel} backwards.} Writes the wrong register, silently.
\item \textbf{\code{SPECIAL} decoded on opcode alone.} Then every R-type is
\code{ADD}.
\item \textbf{\code{ADDI} and \code{LUI} given the same \code{IR\_Xsel}.} They
differ on exactly that one signal.
\item \textbf{Raw bit patterns instead of the package constants.}
\end{itemize}

\vspace{10pt}
\hr
{\small Next week: loads, stores, branches and jumps --- four more rows, and the
control unit is finished. The week after, you run the program.}

\end{document}
